\section{Methodology}
\label{sec:method}

\textbf{Testbed.} All experiments ran on one 4-core virtual machine with
\SI{15}{GB} of memory, Docker~29 and Compose. The system under test was pinned
to two cores and the load generator to the other two. Runs in which the
generator used more than 90\% of its cores or ran out of local ports were
discarded (none in the results reported here unless stated).

\textbf{Data.} \num{200000} synthetic students with six subjects each,
generated deterministically from a seed so that the expected answer for any
roll number can be recomputed. 1\% of students receive a revised result
(re-totalling) when corrections are published. No real student data was used.

\textbf{Workload.} Students arrive as an open-loop Poisson
process~\cite{schroeder2006open} whose rate ramps up during \SI{20}{s} before
release (to 20\% of the peak), jumps to the peak within \SI{3}{s} of release,
and decays exponentially with a time constant of \SI{30}{s} over a
\SI{90}{s} window. This is a time-compressed version of a real result day, which
decays over tens of minutes. Students who arrive early reload every
5--\SI{15}{s} until release. Each student has a patience drawn from
5--\SI{20}{s}; after an error the student refreshes with probability 0.8, after
running out of patience with probability 0.9, after a pause of 1--\SI{5}{s},
for at most 20 attempts. 10\% of students look again after seeing the result;
1\% mistype their roll number. These behavioural parameters are assumptions;
Section~\ref{sec:refresh} varies them. Roll numbers are sampled without
replacement (every arriving student is a different student), or from a Zipf
law ($s{=}1.2$) to model hot keys.

\textbf{Load levels.} We first calibrated design~A: the highest constant
arrival rate at which 99\% of students received a correct result within
\SI{2}{s}, found by stepping and bisection, was \CapA\ students/s. Peak
arrival rates are expressed as multiples of this capacity (1$\times$ to
10$\times$). For reference, the JEE Advanced 2016 incident corresponds to
about 5$\times$~\cite{pagalguy2016jee}.

\textbf{Oracle and metrics.} Every response is classified as correct, stale
(the old version after corrections), wrong, missing (404 for an existing
student), soft failure (HTTP 200 without a result), gate or queue response, or
an availability failure (5xx, timeout, connection error). The main outcome is
the share of students who receive a correct result within 2, 5 and
\SI{10}{s} of their first request after release, and at all
(\emph{served}). Cost is the CPU-seconds used by all containers per
\num{1000} students served. Amplification is requests sent per student.

\textbf{Repetitions.} Each configuration was run 3 times with different seeds
in randomised order; we report medians with 95\% percentile-bootstrap
intervals.

\textbf{Generator validation.} At a constant \num{150} requests/s against
design~A, our generator with evenly spaced arrivals and k6's constant-arrival-rate
executor~\cite{k6docs} measured a median latency of \SI{5.9}{ms} and
\SI{5.9}{ms}, and a 99th percentile of 11.7--\SI{12.1}{ms} and
12.7--\SI{14.0}{ms}, respectively (two runs each).
